The setup and assembly for this experiment was more simple than the math required to make it work. A laser was placed at one end of the table opposite a stepper motor fashioned to a track affixed to an LDR with a diffraction grating on a stand in-between the laser stand and the motor. For the purpose of meeting the short deadline of this experiment, the required length of the track was determined ahead of time by taking twice opposite side length to the theoretical angle for the first order peak in air based on equation (1) as the fringes were expected to get closer together as the index increased. A cuvette was placed between the diffraction grating and the LDR. Initially it was placed roughly 1-2mm from the grating, however, the refractive indices produced by equation (9) carried a 10-20\% error and it was determined that any distance travelled by the light between the grating and the LDR greater than 1mm needed to be accounted for. Thus, the cuvette was placed flush with the grating for equation (8) to be sufficient. Finally, the black box was added to isolate the light reactive portions of the assembly to avoid external interference to maximize the difference between bright and dark fringes.

Figure 2: Diagram of the setup for the experiment including required connections and approximate beam path.

After the code was finished, a measurement was taken of the fringes with the only medium being air to ensure the experimental setup and formulae were able to determine a correct value for the refractive index of air, 1. As a secondary verification, it is possible to place a thin glass slide of thickness $p$ flush with the diffraction grating to ensure the desired spread of the fringes is observed. Next a measurement was taken with the empty cuvette flush with the grating to ensure the estimate that the contribution due to the parallel quartz walls of the vessel was minimal and determine a baseline against which to measure the change in position of the fringes. The cuvette was then filled with each of the available substances one by one ranging in refractive indices and rinsed with ethanol in between each to ensure to minimize contamination as only one of each cuvette was available. The use of a macro cuvette is recommended as initially a plastic semi-micro cuvette was used however, the edges of the cuvette tended to clip the first order fringes whereas, in further attempts, a quartz macro cuvette was used and significantly eased the alignment and data collection processes.

The liquids were measured in order from lowest theoretical refractive index to highest with a consideration for viscosity. For each liquid, the intensity was plotted to ensure a visible peak could be established at the areas of constructive interference. The centers of these peaks were used to determine the displacement of the first order positive and negative peaks and were placed into equation (9) to calculate the refractive indices of each substance. The first laser used has an approximate wavelength of 650nm, to verify the precision of this method a second laser with a 532nm wavelength was exchanged for the 650nm. Twenty readings were taken by the LDR every 0.2mm and averaged to attain the maximum precision possible with the instrument. The inverse of these values were plotted against the position of the LDR on the track centered about the center fringe. The bright spots on the graph were thick, messy peaks with an uncentered maxima, presumably due to limitations of the LDR, so the center of each peak was calculated based on the span of the position that reached a threshold intensity value, and these points were used to determine the distances from the center fringe. The position values were compared to the baseline values established with the empty cuvette fringe positions and plugged into equation (9) to find the respective refractive indices of the fluids.

